% Generated draft: requires mathematical and source review.
\documentclass[11pt]{article}
\usepackage[margin=25mm]{geometry}
\usepackage{fontspec}
\setmainfont{Latin Modern Roman}
\usepackage{amsmath,amssymb,mathtools,unicode-math}
\setmathfont{Latin Modern Math}
\usepackage{enumitem,needspace,xurl,hyperref}
\setlength{\parindent}{0pt}
\setlength{\parskip}{5pt}
\setlength{\emergencystretch}{3em}
\newcommand{\N}{\mathbb{N}}
\newcommand{\Z}{\mathbb{Z}}
\newcommand{\Q}{\mathbb{Q}}
\newcommand{\R}{\mathbb{R}}
\newcommand{\C}{\mathbb{C}}
\newcommand{\F}{\mathbb{F}}
\newcommand{\A}{\mathbb{A}}
\newcommand{\PP}{\mathbb{P}}
\newcommand{\E}{\mathbb{E}}
\newcommand{\eps}{\varepsilon}
\newcommand{\dd}{\,\mathrm{d}}
\begin{document}
\section*{UnsolvedMath ID: 30006589}
\textbf{Source problem code:} catalogue entry
\subsection{Definitions and mathematical statement}
Let
\[
H^{\geq 2}=\Q\langle z_2,z_3,\ldots\rangle
\]
be the free associative algebra over \(\Q\) on the indicated generators. The source considers a map \(G\) sending words to multiple Eisenstein series. To interpret its kernel, take \(G\) to mean the intended \(\Q\)-linear extension, with target the \(\Q\)-vector space \(\mathcal E\) spanned by its values.

The supplied record uses a product \(*\) on \(H^{\geq 2}\), a relation map \(R\colon H^{\geq 2}\times H^{\geq 2}\to H^{\geq 2}\), and an operator \(\theta\colon H^{\geq 2}\to H^{\geq 2}\), but does not give their definitions. With these symbols interpreted as the specific operations intended in the cited report, set
\[
\operatorname{DR}^{*}
:=\operatorname{span}_{\Q}
\left\{\theta^n\bigl(R(u,v)\bigr)*w:
 n\in\Z_{\geq 0},\quad u,v,w\in H^{\geq 2}\right\},
\qquad \theta^0=\operatorname{id}.
\]
Here the span consists of finite rational linear combinations. The conjecture is
\[
\operatorname{DR}^{*}=\ker G
=\{h\in H^{\geq 2}:G(h)=0\}.
\]
Equivalently, for every \(h\in H^{\geq 2}\), the condition \(G(h)=0\) holds if and only if there exist \(m\in\Z_{\geq 0}\), coefficients \(a_i\in\Q\), integers \(n_i\geq 0\), and elements \(u_i,v_i,w_i\in H^{\geq 2}\) such that
\[
h=\sum_{i=1}^{m}a_i\,\theta^{n_i}\bigl(R(u_i,v_i)\bigr)*w_i.
\]

\textit{Unresolved specifications.} This is a statement relative to the source's intended operations, not a self-contained definition of them. The supplied record does not specify the formulas for \(*\), \(R\), or \(\theta\), the normalization and analytic domain of the multiple Eisenstein series, or the treatment of the empty word under \(G\). These data must be recovered from the cited source before the conjecture can be made fully explicit.

For context, the excerpt introduces the intended formal quotient \(\mathcal E^{\mathrm{ff}}=H^{\geq 2}_{*}/\operatorname{DR}^{*}\) and says that the conjecture would identify it with \(\mathcal E\). It also attributes a conjectural \(\mathfrak{sl}_2\)-algebra structure on \(\mathcal E\) to its reference [1], and reports such a structure for the formal quotient, with a proof comparing Fourier expansions. The operations defining these structures and the identities of the excerpt's references [1] and [2] are not supplied.
\subsection{Short English statement}
The conjecture says that every rational relation among the intended multiple Eisenstein series is a finite rational combination of derivatives of the basic relations \(R(u,v)\), multiplied using \(*\) by arbitrary elements of \(H^{\geq 2}\), and that every such combination is a relation. The supplied record leaves the defining operations and analytic conventions unspecified, so this statement requires those source definitions and human mathematical review.
\subsection{Sources}
\begin{itemize}
\item Local UnsolvedMath registry: supplied record with permanent ID \(30006589\), title \emph{Kernel Presentation for Multiple Eisenstein Relations}, exact source problem code \texttt{catalogue entry}, in the Oberwolfach Reports Open Problems collection. No individual proposer is supplied; the proposed year is 2026.
\item \emph{Multiple Zeta Values and Modular Forms} (2026), Oberwolfach Reports, report 7/2026, pp. 433--486, DOI \texttt{10.4171/owr/2026/7}. Supplied source and literature link: \url{https://doi.org/10.4171/owr/2026/7}. The record describes these reports as documenting workshops and their open-problem discussions.
\item Bachmann--Kanno, \emph{Relations and Derivatives of Multiple Eisenstein Series}. Supplied literature link: \url{https://arxiv.org/abs/2602.08176}.
\item Supplied dated literature curation, checked \texttt{2026-08-22}: the record labels the problem open and describes the equality \(\operatorname{DR}^{*}=\ker G\) as conjectural. It states that the cited 2026 work proves an \(\mathfrak{sl}_2\) structure for the formal quotient but does not establish equality with the analytic kernel. This is supplied dated curation derived from the OWR record, not an independently verified literature assessment; its time-sensitive status and claims require independent verification.
\item The registry labels the statement \texttt{corrected\_verified} and says it was reconstructed and checked against the official report. These are supplied verification claims, not verification performed for this draft. The excerpt attributes the main conjecture to reference [2] and the analytic \(\mathfrak{sl}_2\)-algebra conjecture to reference [1], without supplying their numbered bibliography entries.
\end{itemize}
\end{document}
